\section{RWM-U: reuse offline robot data with uncertainty penalties}\label{sec:rwmu}
\textbf{Li, Krause, and Hutter, 2025/2026.} Uncertainty-Aware Robotic World
Model Makes Offline Model-Based Reinforcement Learning Work on Real Robots~\cite{rwmu}.
\loc{Version 3, January 8, 2026; Sections 3.2 and 4; Equations 1--5}

\subsection{Reported mechanism in common notation}
RWM-U learns autoregressive observation dynamics from offline data. Gaussian
prediction heads share recurrent features; disagreement estimates epistemic
uncertainty. MOPO-PPO trains a policy on imagined rollouts with an uncertainty
penalty. Its policy estimator is PPO, not BPTT through model Jacobians. The
paper evaluates robotic manipulation and locomotion, including hardware.
The latest reviewed version is a substantial revision of the 2025 preprint;
the source version is important when checking claims.

For a history $\Hhist_t$ containing observed or recursively predicted inputs,
\begin{align}
m_t&=f_\phi(\Hhist_t),\qquad
\hat o_{t+1}^{(b)}\sim\mathcal N\big(\mu_{\phi_b}(m_t,a_t),\Sigma_{\phi_b}(m_t,a_t)\big),\\
L_{\rm WM}&=\frac1{BK}\sum_b\sum_{k=1}^Kw_k
\E_{\mathcal D}\left[-\log p_{\phi_b}(o_{t+k}\mid\widehat\Hhist_{t+k-1},a_{t+k-1})\right],\\
\bar o_{t+1}&=\frac1B\sum_b\mu_b,\qquad
\Sigma_{U,t+1}=\operatorname{Cov}_b(\mu_b),\\
\tilde r_t&=r_t-\beta_U U_t.
\end{align}
$U_t$ denotes the scalar uncertainty score used by policy optimization; the
paper's variance notation requires a scalar reduction in a practical reward.
Its scaling and the observation normalization should be checked in code before
reproduction. We do not invent an unreported component weighting here.

Writing $\rho_t(\theta)=\policy(a_t\mid o_t)/\pi_{\theta_{\rm old}}(a_t\mid o_t)$,
the normalized clipped PPO objective is
\begin{equation}\label{eq:ppo}
\max_\theta\E\left[\min\left(\rho_t(\theta)\widehat A_t,
\operatorname{clip}(\rho_t(\theta),1-\epsilon,1+\epsilon)\widehat A_t\right)\right],
\end{equation}
where advantage targets use the penalized imagined reward. This explains the
optimization route without reproducing all implementation coefficients.

\subsection{The connection to pretrained-model reuse}
An offline controller cannot correct its model by collecting new interaction
while it improves. A penalty can discourage unsupported imagined behavior, but
it cannot reveal a surface absent from both the data and inputs. If stair contact
is not represented, the conservative solution may be to avoid the stair rather
than learn to traverse it.

For our question, this suggests separating two regimes. First, how far can a
fixed model and fixed dataset support a new controller? Second, once additional
contact data is allowed, which observations most efficiently expand that support?
These should be measured as different experiments. A zero-additional-data
result within an existing dataset is not equivalent to adaptation to a genuinely
new contact distribution.

Shared recurrent features can also induce correlated errors across prediction
heads. Small disagreement does not rule out a common representation failure.
Use independently initialized or independently trained models as a diagnostic
comparison, and measure whether uncertainty predicts errors near new contacts.
The purpose is to test calibration and coverage, not to assume that more ensemble
members automatically provide a conservative bound.

\subsection{A useful comparison with gradient-based methods}
Use the same adapted dynamics and data with PPO in imagination and with pathwise
actor optimization. This separates the value of more imagined data from the
value of transition derivatives. Match the reward penalty and action coordinates;
otherwise a difference may come from pessimism or exploration rather than the
gradient estimator.

An uncertainty penalty, an adaptive rollout horizon, and a request for new
simulator data are three different responses to model ignorance. Their relative
benefit is a research variable. RWM-U motivates the offline control baseline;
it does not supply a measured data budget for our proposed source-to-target
geometry shift.

\section{SuperMimic: online dynamics shared across a motion library}\label{sec:supermimic}
\textbf{Liao et al., 2026.} SuperMimic: Scaling Humanoid Whole-Body Control
with One Differentiable World Model~\cite{supermimic}.
\loc{Sections 3.1--3.3, Equations 1--3; Appendices A--B and D}

\subsection{Reported mechanism in common notation}
SuperMimic alternates simulator collection, autoregressive world-model fitting,
and imagined policy/value updates. The deterministic model conditions on known
simulation context $\xi$; sensor noise is added separately. The policy is goal
conditioned. A motion encoder and shared decoder provide the control interface,
while a critic bootstraps short imagined returns. Prediction heads can be
learned or analytic. The original model weights $\omega$ become $\phi$ here.

The reported general model interface becomes
\begin{align}
\hat x_{t+1}&=\Fhat(\hat x_t,a_t,\xi),\\
\hat r_t&=\widehat R_\phi(\hat x_t,a_t,\Gtask_t),\qquad
\bar{\hat o}_t=O_\phi(\hat x_t,\xi),\\
\hat d_{t+1}&=D_\phi(\hat x_{t+1},\Gtask_t),\qquad
a_t=\policy(\bar{\hat o}_t+\eta_t,g_t).
\end{align}
The reward index is shifted to our transition convention. With $K=8$ and logged
executed actions, a normalized form of the fitting loss is
\begin{equation}
L_{\rm WM}=\E_{\mathcal D}\sum_{k=1}^{K}\left[
\norm{\hat x_k-x_k}^2+(\hat r_{k-1}-r_{k-1})^2
+\norm{\bar{\hat o}_k-\bar o_k}^2+\operatorname{BCE}(\hat d_k,d_k)\right].
\end{equation}
Analytic heads have no trainable fitting term. State differences and scaling
must be adapted to the chosen representation rather than interpreted as a
universal quaternion-coordinate prescription.

The actor uses $H=32$ and soft continuation,
\begin{align}
C_0&=1,\quad C_{t+1}=C_t(1-\hat d_{t+1}),\\
\max_\theta\widehat J_H&=\E\left[\sum_{t<H}\gamma^tC_t\hat r_t
+\gamma^HC_HV_\psi(\hat x_H,g_H,\xi)\right].
\end{align}
Value targets are imagined lambda returns. Model parameters stay fixed during
the actor backward pass, while their input Jacobians remain active.

\subsection{What motion pretraining potentially buys}
Many motion goals can share a physical transition function. This is the
amortization opportunity: changing the desired motion changes actions and rewards,
not the underlying fixed-scene mechanics. A source dataset containing diverse
velocities, configurations, falls, and recoveries may therefore initialize a
useful portion of target dynamics learning.

That statement is a hypothesis about coverage, not an explicit factorization of
network weights. A state MLP can mix actuator dynamics, flat-plane contact, and
source-scene correlations. Neither motion diversity nor long rollout stability
establishes correct behavior when a foot encounters an unseen vertical riser.

Context $\xi$ is also not the same as spatial geometry. Friction, calibration,
and center-of-mass variation can be supplied as episode parameters without telling
the model where a step is. A model that conditions on physical randomization
does not thereby become a terrain-aware model.

\subsection{The actor interface is a second transfer problem}
The normalized encoder--decoder form is
\begin{equation}
\zeta_t=E_\theta(g_t),\qquad
a_t=D_\theta(\Hhist_t,\zeta_t).
\end{equation}
Here $\zeta_t$ is a goal embedding, while $u_t=\mathcal U(a_t)$ continues to mean
the physical actuator command from Section~\ref{sec:notation}.
An extension to new terrain might give the decoder an additional $e_t$ input.
Whether to preserve, adapt, or expand the old goal representation is then an
ablation, not something solved merely by fine-tuning $\phi$.

Reference-consistent rewards matter too. A rigidly timed flat-ground reference
may oppose valid stair motion, particularly root-height change or contact timing.
One can use terrain-compatible references, or retain motion tracking as a soft
prior while optimizing a new task reward. These alternatives should be compared
with the same transition data.

\subsection{Distinguish the paper from our implementation}
Our local project is an implementation of the central training idea, not the
authors' code. It uses a fixed-physics state model with analytic task calculations,
and its world model has no terrain input. The paper's general context-conditioned
formulation is broader than that local baseline. Our recorded short-prefix
evaluations also do not reproduce a full-motion benchmark. This review does not
use those local scores as evidence of published-method parity.

\section{\texorpdfstring{$\nabla$Mimic}{NablaMimic}: stabilize simulator-derived policy updates}\label{sec:nablamimic}
\textbf{Lu and Chen, 2026.} $\nabla$Mimic: Ultra-Efficient Humanoid Control
with Differentiable Simulation~\cite{nablamimic}.
\loc{Sections 4--6; Algorithm 1, Equations 4--5, and Section 8 limitations}

\subsection{Reported mechanism in common notation}
$\nabla$Mimic combines simulator adjoints with median-relative per-sample action
cotangent clipping and repeated regression onto gradient-improved action targets.
NablaEngine specializes the forward and adjoint computation to the G1. Training
modifies problematic geometric derivatives and actuator sensitivities. The
published engine supports flat ground; new terrain requires additional collision
support. The method does not train a world model. Its main recipe uses a
64-step window; the source clipping factor is $\kappa=3$.

The forward window uses $x_{t+1}=\Fsim(x_t,a_t)$ and a value tail. To compare
updates in physical units, let $u_i$ denote the commanded joint setpoint and
$b_i=\partial J_H/\partial u_i$ its action cotangent. The normalized clipping and
target construction is
\begin{align}
\nu_b&=\operatorname{median}_j\norm{b_j},\\
\tilde b_i&=b_i\min\left(1,\frac{\kappa\nu_b}{\norm{b_i}}\right),\\
s_b&=\left(\frac1{Nd_a}\sum_i\norm{\tilde b_i}^2\right)^{1/2},\\
u_i^{\rm target}&=\sg\left(u_i+\delta\,\frac{\tilde b_i}{s_b}\right),\\
L_{\rm regress}(\theta)&=\frac1N\sum_i
\norm{\pi_\theta^u(\Hhist_i,g_i)-u_i^{\rm target}}^2,
\quad\pi_\theta^u=\mathcal U\circ\policy.
\end{align}
Zero norms require the usual numerical convention. The source algorithm clips
inside backward propagation and again on window action gradients; clipping only
the final parameter gradient is not the same operation.

\subsection{Why the update parameterization matters}
Repeated regression reuses a window's gradient information without repeatedly
re-evaluating an increasingly stale physical derivative. Targets stop pulling
once the policy reaches them. The target step $\delta$ is a physical setpoint
scale, rather than an unconstrained scalar multiplier of a potentially enormous
gradient.

For one batch using a common positive median and $\kappa\ge1$, at least half its
clipped action-gradient norms are at least $\nu_b$. Thus
\begin{equation}
s_b\ge\frac{\nu_b}{\sqrt{2d_a}},\qquad
\frac{\norm{u_i^{\rm target}-u_i}}{\sqrt{d_a}}
\le\delta\kappa\sqrt2.
\end{equation}
This normalized bound assumes that common clipping scale; it should not be
blindly applied across batches using different medians. It bounds targets, not
the final nonlinear policy change after arbitrary optimizer steps.

Rescaling coordinates changes gradient norms. If a normalized action maps to
$u=Da+u_0$, then $\nabla_aJ=D^T\nabla_uJ$. Comparing raw action-gradient norms
across controllers with different joint scales can therefore be meaningless.
For our research, evaluate updates in common physical command units or matched
physical effects.

\subsection{What exactness does and does not mean}
The relevant reference is the selected discrete simulator, with contact decisions
held locally fixed where differentiated. It is not an exact derivative of a
physical robot, and the practical training rules intentionally alter some
sensitivities. Clipping also biases the update. These choices may improve
control despite reducing formal derivative fidelity.

The research implication is to measure useful update directions at the scale
actually applied, not simply declare simulator gradients superior or learned
gradients smoother. A terrain adaptation could improve a model's local derivative
while making finite action changes across a contact boundary worse.

\subsection{The open comparison for new geometry}
Use the same target geometry, rewards, policy interface, and reset distribution
when comparing simulator derivatives and model derivatives. First verify that
the simulator supports the required contacts. Otherwise, an apparent failure of
a gradient method is a forward collision-model limitation.

For new stairs, clipping may manage extreme samples but cannot supply missing
terrain observations or discover which new supports exist. Geometry, data
coverage, and gradient conditioning remain separate requirements. The paper
motivates a strong gradient-handling baseline; it does not settle the proposed
motion-pretraining adaptation experiment.

\section{WM-LOCO: predictive terrain context with a PPO controller}\label{sec:wmloco}
\textbf{Liu et al., 2026.} World-Model-Augmented Visual Locomotion for
Humanoids on Foothold-Constrained Terrain~\cite{wmloco}.
\loc{Version 2, November 3, 2026; Section II, Equations 1--11; Sections III--V}

\subsection{Reported mechanism in common notation}
WM-LOCO jointly trains an RSSM from depth, proprioception, actions, and rewards,
then feeds predictive recurrent features to a PPO actor--critic. Its prior
supports action-conditioned rollouts, but the policy estimator is PPO rather
than policy BPTT through those rollouts. Version 2 evaluates stairs, gaps,
stepping stones, and balance beams in simulation and on G1 hardware. It starts
joint training without world-model pretraining.

The paper's deterministic memory becomes $m$, discrete stochastic state becomes $\ell$,
and commanded velocity becomes part of $g_t$. To keep the clocks explicit,
let $t$ index 50 Hz control and $n$ index 10 Hz model updates at $\tau_n=5n$.
Let $\mathbf a_{n-1}$ denote the executed-action information between model updates:
\begin{align}
f_n^o&=\mathcal E_\phi(o_{\tau_n}^{\rm prop},o_{\tau_n}^{\rm depth}),\\
m_n&=f_\phi(m_{n-1},\ell_{n-1},\mathbf a_{n-1}),\\
\ell_n&\sim\mathsf q_\phi(\cdot\mid m_n,f_n^o),\quad
\hat\ell_n\sim\mathsf p_\phi(\cdot\mid m_n),\\
(\hat o_{\tau_n}^{\rm prop},\hat o_{\tau_n}^{\rm depth},\hat r_{\tau_n})&=D_\phi(m_n,\ell_n).
\end{align}
The representation of the intervening action information is an implementation
detail, not an assumption that one control action lasts 100 ms.
Let $\mathsf Q_n$ and $\mathsf P_n$ denote the posterior and prior distributions
from the preceding display. The normalized model objective includes
\begin{align}
L_{\rm WM}&=\E\sum_n\left[L_{\rm prop}+L_{\rm depth}+L_r+L_{\rm KL}\right],\\
L_{\rm KL}&=\beta_d\max\big(\KL(\sg(\mathsf Q_n)\|\mathsf P_n),1\big)
+\beta_r\max\big(\KL(\mathsf Q_n\|\sg(\mathsf P_n)),1\big).
\end{align}
Proprioception reconstruction uses a symlog transform. A policy projection
of the most recent memory supplies
\begin{equation}
f_t^{\rm pred}=\mathcal P_\theta(m_{\lfloor t/5\rfloor}),\qquad
a_t\sim\policy(\cdot\mid o_t^{\rm prop},\Hhist_t,g_t,f_t^{\rm pred}).
\end{equation}
A mixture-of-experts actor--critic and the clipped estimator in
Equation~\eqref{eq:ppo} complete the control path. The model updates at 10 Hz;
the actor runs at 50 Hz and reuses the latest predictive feature between updates.

\subsection{Prediction and policy training are different claims}
Being able to propagate the prior for several steps demonstrates a generative
dynamics capability. Feeding its memory to PPO can improve actions without
differentiating future policy return through that dynamics. Consequently,
success on terrain does not show that the transition Jacobians are adequate
for policy BPTT. The recurrent representation may be useful even when those
Jacobians are poor.

Likewise, successful posterior reconstruction and open-loop prior prediction
are different tests. For a terrain-conditioned policy, evaluate whether memory
preserves support geometry as the viewpoint changes and whether unseen edges
remain distinguishable after propagation. This connects the representation
question to Dreamer's posterior/prior distinction without assuming identical
training algorithms.

\subsection{Task design contributes to terrain success}
The source includes task-specific contact-region rewards alongside locomotion
terms. Their normalized structure can be written
\begin{equation}
r_t=r_t^{\rm loco}-w_b\mathcal C^{\rm boundary}(x_t,E)-w_r\mathcal C^{\rm riser}(x_t,E)
+w_f R_t^{\rm support}.
\end{equation}
The symbols here summarize reward categories, not a replacement for the source's
terrain-specific gating and contact-region definitions. This is important for
comparison: a generic tracking reward and a carefully shaped foothold reward
do not define the same learning problem.

To isolate a world model's contribution, hold those rewards fixed while removing
predictive features, replacing them with a similarly sized non-predictive
encoder, or changing the model pretraining. Otherwise, a better reward or larger
recurrent controller can masquerade as a better dynamics model.

\subsection{Its implication for our research question}
Terrain-aware humanoid world models are already an active research direction.
Simply adding depth or an RSSM to a stair controller is therefore an insufficient
novelty claim. The narrower question is whether motion pretraining reduces new
interaction, and whether model structure improves the utility of its gradients.

WM-LOCO provides a strong predictive-feature baseline for that question. Give it
the same target observations and task specification, and compare data-to-success
with imagined BPTT and decoupled updates. This separates perception, recurrence,
and control-gradient choices instead of conflating them in one success number.
